% ===== tr-08 · b14 · tr08c：第 5 章 p31--33（原文页 31--33）=====
% [原文 p.31]
\begin{proposition}{5.5.4}[NA MA 与实 MA 的比较]
设 $(\mathcal X,\mathcal L)$ 是 $(X_K,L)$ 的半稳定 snc 模型，$\mathrm{Int}(\Delta_J)$ 是 $\Delta_{\mathcal X}$ 的一个 $n$ 维开面。回忆收缩映射 $r_{\mathcal X}:X_K^{an}\to\Delta_{\mathcal X}$。设 $\phi\in C^0(X_K^{an})$ 是半正度量 $\|\cdot\|_{\mathcal L}e^{-\phi}$ 的势，并设 $\phi=\phi\circ r_{\mathcal X}$ 在 $r_{\mathcal X}^{-1}(\Delta_J)$ 上成立。则在 $\mathrm{Int}(\Delta_J)$ 上，NA MA 测度的推前
\[
r_{\mathcal X*}MA(\|\cdot\|e^{-\phi})=n!\,\mathrm{MA}_{\R}(\phi|_{\mathrm{Int}(\Delta_J)})
\]
等于凸函数 $\phi|_{\mathrm{Int}(\Delta_J)}$ 的实 MA 测度（相差因子 $n!$）。
\end{proposition}

% [Tutor Note: $r_{\mathcal X*}$ 是推前（pushforward）而非拉回，星号紧贴下标；右端 $\mathrm{MA}_{\R}$ 的下标是花体 R（$\R$ 宏），表示「实 MA 测度」。原文命题名写作 "(NA MA-real MA comparison)"，译文括注照译。$n!$ 这个因子在 (17) 中会被吸收进分母，正是比较性质的定量内容。]

这一比较的严格证明用到相交理论，下面给出的是启发式解释。考虑动机部分所述代数曲线上的 snc 模型 $\mathcal X$，并进一步假设它是半稳定的。回忆我们的启发式词典：$L\to X_K$ 上的度量 $\|\cdot\|$ 应当编码 $L\to X_t$ 上一族 Hermite 度量 $h_t$，使得在混合拓扑下 $h_t^{1/|\log|t||}\to\|\cdot\|^2$，而 $\|\cdot\|$ 的 NA MA 测度应当是 $h|_{X_t}$ 的曲率形式所对应测度的极限。我们现在聚焦于 $n$ 维开面 $\mathrm{Int}(\Delta_J)\subset\Delta_{\mathcal X}\subset\Delta_{\mathcal X}\sqcup X$ 的邻域，其中有局部坐标 $z_0,\dots z_n$ 满足 $\prod_0^n z_i=t$，以及 $x_i=\frac{\log|z_i|}{\log|t|}$。在局部图景中我们把度量与势等同起来，于是 $\|\cdot\|\sim e^{-\phi}$；忽略 $C^0$ 涨落效应后，$h_t^{1/|\log|t||}\sim e^{-2\phi\circ\mathrm{Log}_{\mathcal X}}$。加上更多光滑性假设，$h_t$ 的曲率形式近似为
\[
|\log|t||\,dd^c\phi\circ\mathrm{Log}_{\mathcal X}=\frac{-1}{2\pi}\sum_{1\le i,j\le n}\frac{\partial^2\phi}{\partial x_i\partial x_j}dx_i\wedge d\arg(z_j).
\]
NA MA 测度应当与极限推前测度
\[
\lim_{t\to 0}\mathrm{Log}_{\mathcal X*}\big((|\log|t||\,dd^c\phi\circ\mathrm{Log}_{\mathcal X})^n\big)=n!\det(D^2\phi)|dx_1\dots dx_n|=n!\,\mathrm{MA}_{\R}(\phi)
\]
一致，后者等于实 MA 测度（相差因子 $n!$）。

% [Tutor Note: 本页 $h_t^{1/|\log|t||}\to\|\cdot\|^2$ 右端是范数平方 $\|\cdot\|^2$（页图可见上标 2），不是 $\|\cdot\|_2$；相应地 $h_t^{1/|\log|t||}\sim e^{-2\phi\circ\mathrm{Log}_{\mathcal X}}$ 里的因子 2 就是从平方来的。$dd^c$ 的 $c$ 是上标，系数原文排作 $\frac{-1}{2\pi}$（分子为 $-1$），不要抄成 $-\frac{1}{2\pi}$ 之外的形式；$\mathrm{Log}$ 是大写对数映射，$\arg$ 取小写正体。]

\begin{remark}{5.12}
在这个启发式计算中，$\phi$ 经由收缩映射分解这一假设使我们能把混合空间 $X\sqcup X_K^{an}$ 替换为其有限逼近 $X\sqcup\Delta_{\mathcal X}$。
\end{remark}

% [Tutor Note: 注 12 的 "factor through the retraction map" 指 $\phi$ 只通过 $r_{\mathcal X}$ 依赖 $X_K^{an}$ 上的点（即 $\phi=\phi\circ r_{\mathcal X}$ 型），故可用有限逼近 $X\sqcup\Delta_{\mathcal X}$ 代替整个混合空间；这正是命题 5.4 假设的用处。]

\subsection{NA Calabi 猜想}
NA 多复变势理论的核心结果是解出了 Calabi 猜想的 NA 类比。好的综述见 [5]。

\begin{theorem}{5.5.5}
[4] 设 $X_K$ 是由代数退化族作基变换得到的光滑射影 $K$-概形。设 $L$ 是 $X_K$ 上的丰沛线丛，$d\mu$ 是支撑在 $X_K$ 的某个 snc 模型的对偶相交复形上的 Radon 概率测度。则 $L$ 上存在唯一的连续半正度量 $\|\cdot\|$，使得
\[
MA(\|\cdot\|)=(L^n)d\mu.
\]
\end{theorem}

% [Tutor Note: 定理 5.5 在原文中整段用斜体排印（与本页其他定理不同，属原文排版习惯），译文按中文惯例不设斜体。引用 [4] 在原文中紧跟在 "Theorem 5.5." 之后，不是定理名，故放在正文开头而不放进可选参数。$(L^n)$ 是相交数，$d\mu$ 是概率测度，故 $\mathrm{MA}$ 的总质量被正规化为 $(L^n)$。]

% [Tutor Note: 定理 5.5 的关键强假设是「$d\mu$ 支撑在某个 snc 模型的对偶相交复形上」：测度并不定义在整个 $X_K^{an}$ 上，而只落在 $\Delta_{\mathcal X}\subset X_K^{an}$ 里，正因如此才能用相交理论把 $\mathrm{MA}(\|\cdot\|)$ 与 $(L^n)d\mu$ 比较；解的唯一性同样依赖这一归一化（总质量被固定为 $(L^n)$）。]

% [原文 p.32]
他们的策略采用「变分方法」。在 $L$ 上连续半正度量的空间（等价地看作连续 $\theta$-psh 势 $\phi$）上存在一个凹能量泛函 $\mathcal E$，其第一变分由 NA MA 测度给出。人们寻求泛函
\[
F_\mu(\phi)=\mathcal E(\phi)-(L^n)\int_{X_K^{an}}\phi\,d\mu
\]
的最大值点：先把 $\phi$ 的空间扩大为函数空间 $PSH(X_K^{an},\theta)$，它模去加实数常数后是紧的；这类似于 Kähler 情形下 $PSH(X,\omega)/\R$ 的 $L^1$-紧性。NA MA 测度与能量泛函的概念自然地延拓到 $PSH(X_K^{an},\theta)$ 内部的能量类函数上，这与复多复变势理论的情形十分类似。然后人们证明该最大值点实际上是一个临界点，即 NA MA 方程的弱解。这一点颇为微妙，因为 $PSH(X_K^{an},\theta)$ 中函数的小扰动可能因失去正性而落到该类之外。人们用 Kolodziej 估计的类比来证明弱解的连续性。解的唯一性再次依赖于 $\mathcal E$ 的凹性。尽管这一策略与复解析情形具有非常相似的逻辑结构，其技术基础却建立在相交理论与双有理几何中的消没定理之上，而非微分算子。

% [Tutor Note: 能量泛函的字母是花体 $\mathcal E$（页图为花体 E），与线丛 $L$、$X$ 的普通斜体必须区分；文字层抽成 "E"，照抄会与别的 $E$ 混淆（本章前文 $E_i$ 是除子）。$F_\mu$ 的下标是测度 $\mu$。原文的 $PSH$ 三个字母（以及 $Sk$、$MA$、$CY$ 等）均为数学斜体，不要排成正体缩写；只有 $\mathrm{Log}$、$\mathrm{Re}$、$\mathrm{Const}$ 等在页图上是正体的才用 \mathrm。]

我们主要关心的是 $X_K$ 来自大复结构极限的情形。此时 NA 多复变势理论给出
\begin{equation}
MA(\|\cdot\|_{CY})=(L^n)d\mu_0, \tag{5.16}
\end{equation}
的唯一解，其中 $d\mu_0$ 是支撑在本质骨架 $Sk(X)\subset X_K^{an}$ 上的 Lebesgue 测度（参看第 5.3.1 节）。我们称 $\|\cdot\|_{CY}$ 为非阿基米德 Calabi--Yau 度量。

% [Tutor Note: (16) 左端的 $MA$ 原文用数学斜体；而 (17) 同一记号原文改用正体 $\mathrm{MA}_{\R}$------这是原文内部的排印不一致，译文按页图分别照抄（其余各处同为斜体）。$\|\cdot\|_{CY}$ 的下标 $CY$ 是正体大写，指 Calabi--Yau。]

\subsection{比较性质}
关于非阿基米德 CY 度量 $\|\cdot\|_{CY}$，除存在性与连续性之外，已知的结果很少。下面我们讨论如下猜想性的 NA MA 与实 MA 比较性质的含义。

\begin{definition}{5.5.6}
我们称 $\|\cdot\|_{CY}$ 满足 NA MA 与实 MA 比较性质，如果存在 $(X,L)$ 的半稳定 snc 模型 $(\mathcal X,\mathcal L)$，使得由 $\|\cdot\|_{CY}=\|\cdot\|_{\mathcal L}e^{-\phi_0}$ 定义的势 $\phi_0$ 在收缩映射于所有 $n$ 维开面 $\mathrm{Int}(\Delta_J)\subset Sk(X)$ 上的原像上满足 $\phi_0=\phi_0\circ r_{\mathcal X}$。
\end{definition}

注意 $\mathrm{Int}(\Delta_J)\subset\Delta_{\mathcal X}$ 继承了一个自然的整仿射结构。由于按命题 5.5.2，$\phi_0$ 的限制在这些面上是凸的，其实 MA 测度有意义；并且由命题 5.5.4，它在 $\mathrm{Int}(\Delta_J)$ 上满足实 MA 方程
\begin{equation}
\mathrm{MA}_{\R}(\phi_0)=\frac{(L^n)}{n!}d\mu_0. \tag{5.17}
\end{equation}

% [Tutor Note: 定义 5.5.6 里的 "it satisfies the real MA equation on Int(Δ_J)" 与 (17) 是同一陈述的两种写法；注意 (17) 中 $n!$ 出现在分母，与命题 5.5.4 的「相差因子 $n!$」正好对应。$Sk(X)$ 在原文中排作两个斜体字母 $Sk$，不是正体 operator。]

% [原文 p.33]
几点评论如下：
\begin{itemize}
\item 比较性质是代数/非阿基米德几何中的一个猜想，先验地（a priori）并不涉及 PDE 概念；它的 PDE 推论是后验地（a posteriori）得到的。
\item Kontsevich--Soibelman 图景（参看第 5.3.3 节）预期：在本质骨架上、在某个奇异轨迹之外存在实 MA 方程的解。而借助比较性质的 NA-MA 方程是已知唯一能系统产生解的方法。
\item 在环面簇上环面不变度量（toric invariant metrics）的语境中，NA MA 测度与实 MA 方程之间已有比较结果，参看 [28, Prop. 4.4.4]。
\item 读者可能觉得，NA 多复变势理论是解实 MA 方程的一条非常冗长的途径。然而出乎意料的是，人们甚至不知道如何在 $Sk(X)$ 上整体地表述实 MA 方程------不仅不知道在 $n$ 维面上如何表述，也不知道在低维面上如何表述。

一个问题是：$Sk(X)$ 并不带有显然的预定仿射结构，而只有分段仿射结构，因此没有显然的、与坐标无关的实 MA 测度定义。看来 Kontsevich 与 Soibelman 所猜想的仿射结构需要与实 MA 方程同时求解，颇像自由边界 PDE 问题。

另一个问题是：解实 MA 方程需要先指定允许作为势的凸函数类，正如解复 Monge--Ampère 方程需要先指定 Kähler 势的含义。我们目前不知道在像 $Sk(X)$ 这样的分段仿射流形上定义凸势类的任何直接途径。Berkovich 空间 $X_K^{an}$ 上的半正度量尽管抽象，却是我们唯一可用的替代品。
\item 若人们相信 NA Calabi--Yau 度量 $\|\cdot\|_{CY}$ 是 $X_t$ 上 Calabi--Yau 度量在混合拓扑下的势论极限，且 $\|\cdot\|_{CY}$ 的信息可以由 $Sk(X)$ 上的数据恢复，那么人们也许会倾向于认为 $\|\cdot\|_{CY}$ 的势整体地经由某个收缩映射 $X_K^{an}\to Sk(X)$ 分解，而该映射或许通过某个因子化对数终端极小模型来定义。

这样的陈述需要面对一个困难：因子化对数终端模型未必唯一，原则上收缩映射依赖于模型的选取。NA MA 方程将如何选出一个预定的收缩映射，看来是高度非平凡的。
\end{itemize}

比较性质的表述比这更谨慎。我们允许对偶相交复形 $\Delta_{\mathcal X}$ 严格大于 $Sk(X)$，并且对 $Sk(X)$ 各 $n$ 维面的补集不作任何假设。

% [Tutor Note: 第 2 条中点出了 "NA-MA equation"（带连字符），与本章其余处的 "NA MA"、"NA MA measure" 写法不同，原文并不统一，译文随原文。末段 "We allow the dual intersection complex $\Delta_{\mathcal X}$ to be strictly bigger than $Sk(X)$" 一句起于 p33 页末、续到 p34 首行 "Sk(X), and there is no assumption..."，由本块（b14）写完整句；tr08d 从 p34 的下一句 "Regarding the problem above..." 开始。]
